\documentclass[11pt]{article}

\usepackage{amsmath,amssymb}
\usepackage{booktabs}
\usepackage{graphicx}
\usepackage{siunitx}
\usepackage[margin=2.5cm]{geometry}
\usepackage{float} % For better figure placement

\sisetup{
  round-mode=places,
  round-precision=6,
  scientific-notation=false
}

\title{Meta-Optimization of a Real-Coded Genetic Algorithm\
for Continuous Benchmark Functions}

\author{Negura Teodor Alexandru\
{\small ``University A.I. Cuza''}}

\date{The Second Assignment -- Extended Study}

\begin{document}

\maketitle

\begin{abstract}
This report extends the work from Homework~2 by introducing a meta-optimization
layer on top of the real-coded Genetic Algorithm (GA) previously designed for
continuous function minimization.
The goal of this extended study is to automatically tune the main GA parameters
(population size, crossover probability, mutation probability and the
distribution indices of SBX and polynomial mutation) via a higher-level
\emph{meta-GA}, and to quantify the performance gain over the fixed-parameter
GA from Homework~2.

The study considers the same four benchmark functions as before---De~Jong~1,
Rastrigin, Schwefel and Michalewicz---in dimensions $d \in \{5,10,30\}$.
For each $(\text{function}, d)$ pair, a dedicated meta-GA searches in the
parameter space of the underlying GA.
The best configuration found by the meta-GA is then re-evaluated in 30
independent runs and compared to the Homework~2 baseline in terms of solution
quality (best, mean, standard deviation) and relative improvement.

The results show that for functions where the baseline GA already reaches
machine precision (De~Jong~1 and low-dimensional Rastrigin), meta-optimization
cannot bring substantial gains.
In contrast, on more challenging landscapes (Schwefel and Michalewicz) the
meta-optimized GA achieves consistent improvements of up to about
$13\%$ in mean fitness, while also reducing variance in most cases.
\end{abstract}

\section{Introduction}

Continuous black-box optimization appears in many scientific and engineering
applications: parameter tuning of physical models, machine learning,
signal processing, control, and many others.
Often the objective function is non-convex, noisy or non-differentiable,
making gradient-based methods inadequate or highly sensitive to initialization.

In this context, metaheuristics such as hill climbing, Simulated Annealing
(SA) and Genetic Algorithms (GAs) provide robust, general-purpose tools for
global optimization.
Homework~1 focused on local-search methods (several hill-climbing variants
and SA), while Homework~2 moved to a population-based approach and designed
a real-coded GA capable of reaching very high precision on four standard
benchmark functions in up to 30 dimensions.

Although the Homework~2 GA already outperformed the local-search baselines
by a large margin, its performance still depended strongly on manually chosen
parameters (population size, crossover and mutation rates, distribution
indices of variation operators).
Choosing these parameters by hand is time-consuming and problem-dependent,
and there is no guarantee that they are close to optimal.

This extended study addresses this limitation by introducing a \emph{meta-GA}:
a second-level genetic algorithm that searches in the parameter space of the
original GA.
The objectives of this report are:
\begin{enumerate}
  \item to define a suitable parameterization and search space for the
        real-coded GA from Homework~2;
  \item to design and implement a meta-GA that automatically tunes these
        parameters on each benchmark function and dimensionality;
  \item to compare quantitatively the performance of the meta-optimized GA
        against the fixed-parameter GA of Homework~2.
\end{enumerate}

\section{Methods}

\subsection{Benchmark functions and problem formulation}

We consider the same four continuous benchmark functions used in Homework~2.

\paragraph{De~Jong~1 (Sphere).}
\begin{equation}
  f_{\text{DJ1}}(x) = \sum_{i=1}^{d} x_i^2,
\end{equation}
with $x_i \in [-5.12, 5.12]$ and global minimum $f^\star = 0$ at $x^\star = 0$.
This is a smooth, strictly convex and unimodal function.

\paragraph{Rastrigin.}
\begin{equation}
  f_{\text{Ras}}(x) = 10d + \sum_{i=1}^{d}
    \bigl( x_i^2 - 10 \cos (2\pi x_i)\bigr),
\end{equation}
with $x_i \in [-5.12, 5.12]$ and global minimum $f^\star = 0$ at $x^\star = 0$.
The function is highly multimodal, with a large number of regularly spaced
local minima.

\paragraph{Schwefel.}
We use the classical Schwefel function with domain $x_i \in [-500,500]$:
\begin{equation}
  f_{\text{Sch}}(x) = - \sum_{i=1}^{d}
    x_i \sin\bigl(\sqrt{|x_i|}\bigr),
\end{equation}
which is deceptive: the global optimum is far from the origin, while many
deep local minima are closer.

\paragraph{Michalewicz.}
\begin{equation}
  f_{\text{Mich}}(x) = -
  \sum_{i=1}^{d} \sin(x_i)
  \left[ \sin\left( \frac{i x_i^2}{\pi} \right) \right]^{2m},
\end{equation}
with $m = 10$ and $x_i \in [0,\pi]$.
The function has sharp, narrow valleys and is difficult in high dimension.

In all cases, the optimization task is
\begin{equation}
  \min_{x \in D \subset \mathbb{R}^d} f(x),
\end{equation}
for $d \in \{5,10,30\}$, as in Homework~2.

\subsection{Baseline real-coded GA from Homework~2}

The baseline algorithm is the real-coded GA designed in Homework~2.
Only a brief summary is given here, since the full details are already
documented.\footnote{% The high-level design and motivation follow standard real-coded GA literature
and were guided by the course notes.}

\paragraph{Representation and initialization.}
Each individual encodes a candidate solution
$x = (x_1,\dots,x_d) \in \mathbb{R}^d$ as a vector of doubles.
This real-coded representation avoids the precision and Hamming-cliff issues
associated with binary encodings.

The initial population is generated by sampling each coordinate independently
from a uniform distribution over the domain of the corresponding function,
e.g.\ $x_i \sim U[-5.12,5.12]$ for De~Jong~1 and Rastrigin.
This ensures that the initial population covers the whole search space and
does not introduce a bias towards any particular region.

\paragraph{Fitness evaluation.}
All four benchmarks are minimization problems.
The GA uses the raw function value as fitness and interprets lower values
as better; no additional rescaling is applied.

\paragraph{Selection: tournament selection.}
Parent selection uses $k$-tournament selection with $k = 3$:
\begin{enumerate}
  \item sample $k$ individuals uniformly at random from the population;
  \item choose the individual with the lowest fitness as the parent.
\end{enumerate}
Tournament selection is robust and independent of absolute fitness values,
making it suitable for functions with negative optima (e.g.\ Michalewicz,
Schwefel).
By tuning $k$, the selection pressure can be controlled; in our experiments
$k = 3$ offered a good balance between exploration and exploitation.

\paragraph{Crossover: Simulated Binary Crossover (SBX).}
For recombination we use Simulated Binary Crossover (SBX), a standard
operator for real-coded GAs:
\begin{itemize}
  \item for each pair of parents $(x^{(1)}, x^{(2)})$ and each gene $i$,
        SBX generates two children $y^{(1)}_i$ and $y^{(2)}_i$ by sampling
        around the parents according to a distribution controlled by a
        parameter $\eta_c$ (the distribution index);
  \item when parents are close in a coordinate, children are also close,
        which favours exploitation; when parents are far apart, SBX explores
        the space between them.
\end{itemize}
We used a crossover probability $p_c = 0.9$ and a distribution index
$\eta_c = 15$.
A high crossover rate, combined with tournament selection and elitism, allows
good building blocks to be propagated through the population.

\paragraph{Mutation: polynomial mutation.}
Mutation uses a polynomial mutation operator, which for each gene perturbs
$x_i$ by a small value $\Delta$ drawn from a polynomial distribution
controlled by a parameter $\eta_m$:
\begin{itemize}
  \item with high probability, mutations are small and refine the current
        solution (exploitation);
  \item with small probability, mutations can be larger jumps, helping the
        algorithm escape local minima (exploration).
\end{itemize}
We used a per-gene mutation probability $p_m = 1/d$ and $\eta_m = 20$,
which biases the search towards fine-tuning in later generations, while
still allowing occasional larger moves in early generations.

\paragraph{Replacement and elitism.}
Generational replacement is used: in each generation a full offspring
population is created and replaces the current one, with \emph{elitism}:
the best individual of generation $t$ is copied unchanged to generation
$t+1$.
This guarantees that the best fitness found so far is monotonically
non-increasing over generations and avoids the loss of good solutions due to
stochastic effects.

\paragraph{Implementation details.}
The GA is implemented in C++ using an object-oriented architecture:
\begin{itemize}
  \item an \texttt{ObjectiveFunctions} module encapsulates the definitions
        and domains of De~Jong~1, Rastrigin, Schwefel and Michalewicz;
  \item a \texttt{GeneticAlgorithm} class manages the population, selection,
        variation, and the main evolutionary loop;
  \item an \texttt{Individual} structure stores the real-valued genotype and
        its fitness.
\end{itemize}
Random numbers are generated using the C++ standard library Mersenne Twister,
with different seeds for each independent run.
The code was instrumented to record, for every run, the best fitness per
generation and the summary of the final best solution (best, worst, mean,
standard deviation, mean error).

\paragraph{Baseline parameter settings.}
In Homework~2, the following population sizes $N$ and number of generations $T$
were used:
\begin{itemize}
  \item $d=5,10$: $N = 150$, $T = 1500$;
  \item $d=30$: $N = 300$, $T = 3000$;
  \item Schwefel in $d=30$ (hardest case): $N = 1000$, $T = 5000$.
\end{itemize}
For each configuration (function, dimension), the GA was run 30 times with
different random seeds and summary statistics of the final best fitness were
reported.

\subsection{Meta-Genetic Algorithm for parameter optimization}

This extended study introduces a second-level GA---the \emph{meta-GA}---which
searches in the parameter space of the baseline GA.
The inner GA (``base GA'') retains its structure (representation, operators),
but its numerical parameters are no longer fixed by hand; instead they are
encoded as genes of the meta-individuals.

\subsubsection*{Parameter vector and search space}

Each meta-individual encodes a parameter vector
\[
  \theta = (N, p_m, p_c, \eta_c, \eta_m),
\]
where:
\begin{itemize}
  \item $N \in \{300,\dots,1000\}$ is the population size (integer);
  \item $p_m \in [0.01,0.40]$ is the per-gene mutation probability;
  \item $p_c \in [0.60,0.95]$ is the crossover probability;
  \item $\eta_c \in [5,40]$ is the SBX distribution index;
  \item $\eta_m \in [5,40]$ is the polynomial mutation distribution index.
\end{itemize}
The bounds are chosen to cover a broad range around the values used in
Homework~2, while avoiding obviously degenerate configurations.

In practice, $N$ is encoded as a real variable and rounded to the closest
integer when instantiating the base GA.

\subsubsection*{Fitness of a meta-individual}

For a fixed benchmark function $f$ and dimension $d$, the fitness of a
meta-individual $\theta$ is obtained by running the base GA with parameters
$\theta$ and recording the final best fitness after a fixed number of
generations $T(d)$ (same $T$ as in Homework~2 for that dimension).
To reduce noise, the meta-fitness is defined as the mean best fitness over
$R=5$ independent base-GA runs with the same $\theta$ but different seeds:
\[
  F(\theta) = \frac{1}{R} \sum_{r=1}^{R}
    \min_{x \in D} f(x; \theta, r).
\]
Since we minimize the benchmark functions, the meta-GA also minimizes
$F(\theta)$.

\subsubsection*{Meta-GA operators and settings}

The meta-GA itself is a real-coded GA with the following settings:
\begin{itemize}
  \item population size $N_{\text{meta}} = 40$;
  \item tournament selection of size $k=2$;
  \item SBX crossover with $p_c^{\text{meta}} = 0.9$ and
        $\eta_c^{\text{meta}} = 10$;
  \item polynomial mutation with per-gene rate
        $p_m^{\text{meta}} = 0.2$ and
        $\eta_m^{\text{meta}} = 15$;
  \item generational replacement with elitism (best meta-individual copied);
  \item number of meta-generations $T_{\text{meta}} = 50$.
\end{itemize}

A separate meta-GA is run for each pair $(\text{function}, d)$.
The best meta-individual $\theta^\star$ at the end of the meta-run is
interpreted as the \emph{optimized parameter configuration} for that
problem.
These configurations are listed in Section~\ref{sec:opt-params}, where they
are compared against the baseline parameters and results from Homework~2.

\section{Experimental setup}

\subsection{Dimensions, repetitions and stopping criteria}

The experiments reuse the same dimensionalities as Homework~2:
$d \in \{5,10,30\}$.
For each function and dimension, we distinguish two phases:

\begin{enumerate}
  \item \textbf{Meta-optimization phase.}
        The meta-GA is run on the chosen $(\text{function},d)$ until
        $T_{\text{meta}}=50$ generations.
        Its output is the best parameter vector
        $\theta^\star = (N, p_m, p_c, \eta_c, \eta_m)$.

  \item \textbf{Evaluation phase.}
        The base GA is run with parameters $\theta^\star$ for
        $T(d)$ generations (same $T$ as in Homework~2) and evaluated over
        30 independent runs.
        The same was done in Homework~2, but with fixed parameters instead
        of $\theta^\star$.
\end{enumerate}

This allows a direct comparison between:
\begin{itemize}
  \item the \emph{baseline GA} (Homework~2) with fixed parameters;
  \item the \emph{meta-optimized GA} (this extended study) with parameters
        tuned by the meta-GA.
\end{itemize}

\subsection{Measured statistics}

For both the baseline and the meta-optimized GA, and for every
(function, dimension) pair, we report:
\begin{itemize}
  \item best (minimum) fitness over 30 runs;
  \item mean and standard deviation of the final best fitness;
  \item worst (maximum) final best fitness;
  \item relative improvement in mean fitness:
        \[
          \text{Improvement} =
          \frac{\text{Baseline mean} - \text{Optimized mean}} 
               {|\text{Baseline mean}|} \times 100\%.
        \]
\end{itemize}

\section{Experimental results}

\subsection{Baseline results from Homework~2}
\label{sec:baseline}

To provide context for the meta-optimization results, we first present
the baseline performance of the fixed-parameter GA from Homework~2.
Table~\ref{tab:baseline} shows the results obtained with the manually
tuned parameters: $p_c = 0.9$, $\eta_c = 15$, $p_m = 1/d$, $\eta_m = 20$,
and the population sizes and generation counts specified in Section~2.2.

\begin{table}[t]
  \centering
  \small
  \caption{Baseline GA results from Homework~2 for each benchmark function
  and dimension (30 runs per configuration). These results serve as the
  baseline for evaluating the improvements achieved by meta-optimization.}
  \label{tab:baseline}
  \begin{tabular}{lrrrrrr}
    \toprule
    Function & $d$ & Best & Mean & Std. dev. & Worst & Mean $\Delta$ \\
    \midrule
    De~Jong~1 & 5  & 0.000000 & 0.000000 & 0.000000 & 0.000000 & 0.000000 \\ 
    De~Jong~1 & 10 & 0.000000 & 0.000000 & 0.000000 & 0.000000 & 0.000000 \\ 
    De~Jong~1 & 30 & 0.000000 & 0.000000 & 0.000000 & 0.000000 & 0.000000 \\ 
    \midrule
    Rastrigin & 5  & 0.000000 & 0.000000 & 0.000000 & 0.000001 & 0.000000 \\ 
    Rastrigin & 10 & 0.000000 & 0.000002 & 0.000002 & 0.000007 & 0.000002 \\ 
    Rastrigin & 30 & 0.000018 & 0.000034 & 0.000012 & 0.000085 & 0.000034 \\ 
    \midrule
    Schwefel & 5  & -2094.914436 & -1925.152823 & 121.106053 & -1621.161098 & 169.761677 \\ 
    Schwefel & 10 & -4189.828868 & -3759.502918 & 187.225829 & -3479.198862 & 430.326082 \\ 
    Schwefel & 30 & -11621.979816 & -11089.007358 & 279.858859 & -10556.034850 & 1480.479642 \\ 
    \midrule
    Michalewicz & 5  & -4.687658 & -4.686266 & 0.007497 & -4.645895 & 0.002006 \\ 
    Michalewicz & 10 & -9.660152 & -9.649157 & 0.031539 & -9.497392 & 0.011094 \\ 
    Michalewicz & 30 & -29.615087 & -29.355069 & 0.302678 & -28.616124 & 0.275831 \\ 
    \bottomrule
  \end{tabular}
\end{table}

As Table~\ref{tab:baseline} shows, the Homework~2 GA already achieved
excellent results on De~Jong~1 and Rastrigin, reaching machine precision
in most cases.
However, on Schwefel and Michalewicz there is still room for improvement,
with noticeable variance across runs and mean fitness values that are not
at the theoretical optimum.

\subsection{Optimized parameter configurations}
\label{sec:opt-params}

Table~\ref{tab:opt-params} summarizes the parameter configurations selected
by the meta-GA for each benchmark function and dimensionality.
For each configuration we report the population size $N$, mutation and
crossover probabilities $(p_m, p_c)$, distribution indices
$(\eta_c,\eta_m)$, and the average fitness used during the meta-optimization
phase.

\begin{table}[t]
  \centering
  \small
  \caption{Meta-optimized GA parameters for each benchmark function and
  dimension.
The ``Avg. fitness'' column shows the mean final best fitness used by
  the meta-GA during parameter search (training criterion).}
  \label{tab:opt-params}
  \begin{tabular}{lrrrrrr}
    \toprule
    Function & $d$ & $N$ &
    $p_m$ & $p_c$ & $\eta_c$ & $\eta_m$ \\
    \midrule
    De~Jong~1 & 5  & 530 & 0.093410 & 0.914757 & 26.856936 & 14.145203 \\ 
    De~Jong~1 & 10 & 662 & 0.056342 & 0.777889 & 26.110097 & 18.119989 \\ 
    De~Jong~1 & 30 & 944 & 0.019342 & 0.863242 & 36.977088 & 14.341666 \\ 
    \midrule
    Rastrigin & 5  & 477 & 0.094621 & 0.869877 & 27.144409 & 33.163233 \\ 
    Rastrigin & 10 & 923 & 0.049221 & 0.942413 & 18.470366 & 16.239478 \\ 
    Rastrigin & 30 & 689 & 0.018929 & 0.838323 & 32.410145 & 37.854956 \\ 
    \midrule
    Schwefel & 5  & 716 & 0.122222 & 0.819668 & 25.855086 &  6.829532 \\ 
    Schwefel & 10 & 817 & 0.355377 & 0.757421 & 23.356420 & 15.202611 \\ 
    Schwefel & 30 & 562 & 0.023579 & 0.759851 & 19.823684 &  5.000000 \\ 
    \midrule
    Michalewicz & 5  & 529 & 0.105795 & 0.679843 & 22.069189 & 12.126143 \\ 
    Michalewicz & 10 & 771 & 0.092374 & 0.795045 & 35.305259 & 24.385418 \\ 
    Michalewicz & 30 & 522 & 0.012661 & 0.615681 & 38.074771 & 22.036382 \\ 
    \bottomrule
  \end{tabular}
\end{table}

Several trends are visible already:
population size tends to increase with dimensionality on De~Jong~1 and
Rastrigin, while for Schwefel the meta-GA prefers a moderately sized
population already in 30D (compared to the very large hand-chosen population
from Homework~2).
Mutation rates are relatively high for 5D and 10D Schwefel and Michalewicz,
encouraging exploration on these rugged landscapes, while in 30D they drop
to small values combined with large distribution indices, favouring fine
tuning near good basins.

\subsection{Performance comparison with the Homework~2 GA}

Table~\ref{tab:comparison} reports the detailed comparison between the
baseline GA (Homework~2, results from Table~\ref{tab:baseline}) and the
meta-optimized GA (this extended study, using parameters from
Table~\ref{tab:opt-params}) on all functions and dimensions.
For each configuration we show the baseline mean fitness, and the best,
mean, standard deviation and relative improvement achieved by the
meta-optimized GA.

\begin{table}[t]
  \centering
  \small
  \caption{Meta-GA vs. Homework~2 baseline GA.
  ``Baseline mean'' refers to the mean fitness from Table~\ref{tab:baseline}.
  ``Optimized'' columns summarize the 30-run performance of the GA with
  meta-optimized parameters from Table~\ref{tab:opt-params}.}
  \label{tab:comparison}
  \begin{tabular}{lrrrrrr}
    \toprule
    Function & $d$ &
    Baseline mean &
    Opt.\ best &
    Opt.\ mean &
    Opt.\ std &
    Improv. (\%) \\
    \midrule
    De~Jong~1 & 5  &  0.000000 &  0.000000 & 0.000000 & 0.000000 &   0.000000 \\ 
    De~Jong~1 & 10 &  0.000000 &  0.000000 & 0.000000 & 0.000000 &   0.000000 \\ 
    De~Jong~1 & 30 &  0.000000 &  0.000000 & 0.000000 & 0.000000 &   0.000000 \\ 
    \midrule
    Rastrigin & 5  &  0.000000 &  0.000000 & 0.000000 & 0.000000 &   0.000000 \\ 
    Rastrigin & 10 &  0.000002 &  0.000000 & 0.000000 & 0.000000 & 100.000000 \\ 
    Rastrigin & 30 &  0.000034 &  0.000000 & 0.000000 & 0.000000 & 100.000000 \\ 
    \midrule
    Schwefel & 5  & -1925.152823 & -2094.914436 & -2094.914436 & 0.000000 &  8.818085 \\ 
    Schwefel & 10 & -3759.502918 & -4188.921208 & -3979.718321 & 132.456766 &  5.857567 \\ 
    Schwefel & 30 & -11089.007358& -12569.486618& -12569.486618& 0.000000 & 13.350873 \\ 
    \midrule
    Michalewicz & 5  & -4.686266 & -4.687658 & -4.687658 & 0.000000 & 0.029708 \\ 
    Michalewicz & 10 & -9.649157 & -9.660152 & -9.659906 & 0.000937 & 0.111398 \\ 
    Michalewicz & 30 & -29.355069& -29.630884& -29.583519& 0.147466 & 0.778229 \\ 
    \bottomrule
  \end{tabular}
\end{table}

\paragraph{De~Jong~1.}
On the unimodal De~Jong~1 function, the Homework~2 GA already reached the
global optimum to machine precision in all dimensions, with zero variance.
Consequently, the meta-GA cannot find any parameter configuration that
improves upon a mean of exactly zero; all statistics remain at numerical zero
for both the baseline and the optimized GA.
This confirms the correctness and robustness of the original implementation.

\paragraph{Rastrigin.}
For Rastrigin, the baseline GA was already extremely accurate, but in 10D
and 30D the mean fitness remained slightly above zero
($2\times10^{-6}$ and $3.4\times10^{-5}$, respectively).
After meta-optimization, the GA consistently finds solutions with mean fitness
numerically equal to zero and zero standard deviation in 30 runs.
In absolute terms, the improvement is tiny, but it shows that the
meta-GA is capable of fine-tuning the parameters enough to completely
eliminate residual numerical error on a highly multimodal landscape.

\paragraph{Schwefel.}
The most significant gains appear on the deceptive Schwefel function.
In 5D, the baseline mean fitness was about $-1925.15$; the meta-optimized GA
improves this to roughly $-2094.91$, an $8.8\%$ improvement in mean fitness.
In 30D, the improvement is even larger:
the baseline mean of $\approx -11089.0$ is improved to
$\\approx -12569.5$, a gain of more than $13\%$.
These results indicate that the large, problem-specific populations and
operator settings discovered by the meta-GA are more appropriate for Schwefel
than the hand-crafted parameters of Homework~2.

\paragraph{Michalewicz.}
On Michalewicz, improvements are more modest but still consistent.
In 5D, the baseline mean of $-4.6863$ is improved very slightly to $-4.6877$.
In 10D and 30D, the improvements are around $0.1\%$ and $0.8\%$ respectively.
While these percentage gains look small, they occur very close to the
(best known) global optimum and thus correspond to non-trivial refinement
of already high-quality solutions.

\subsection{Visual Analysis of Performance and Adaptation}

To further illustrate the performance differences and the complexity of parameter selection, we provide a visual analysis of the experimental results.

Figure~\ref{fig:perf_comparison} presents a direct side-by-side comparison of the mean fitness achieved by the baseline GA versus the meta-optimized GA. This graph explicitly highlights the key finding of this study: while 'solved' problems like DeJong1 show no change, the deceptive Schwefel function sees dramatic improvements (visible as the gap between the red and green bars), validating the efficacy of the meta-optimization approach on difficult landscapes.

\begin{figure}[H]
    \centering
    \includegraphics[width=\linewidth]{graph2_performance_comparison.png}
    \caption{Performance comparison between the Baseline GA (Homework 2) and the Meta-Optimized GA. Note the significant improvement on the Schwefel function across all dimensions.}
    \label{fig:perf_comparison}
\end{figure}

Figure~\ref{fig:correlations} investigates the cost-benefit relationship of meta-optimization. The left panel shows a clear correlation between problem difficulty (baseline error) and the percentage of improvement gained, suggesting that meta-optimization yields the highest return on investment for the hardest problems. The right panel confirms that there is no simple linear relationship between population size and improvement, reinforcing the need for automated tuning.

\begin{figure}[H]
    \centering
    \includegraphics[width=\linewidth]{graph1_correlations_enhanced.png}
    \caption{Cost-Benefit Analysis. Left: Harder problems (larger baseline error) yield greater improvements. Right: No simple correlation exists between population size and performance gain, indicating complex parameter interactions.}
    \label{fig:correlations}
\end{figure}

Finally, Figure~\ref{fig:param_adaptation} demonstrates the adaptive nature of the meta-GA. It reveals that optimal parameters are highly problem-dependent. For instance, while mutation rates typically decrease as dimensionality increases (to preserve good solutions), the deceptive Schwefel 10D case triggered a massive spike in mutation rate (approx 0.35), a strategy the meta-GA automatically discovered to escape local optima.

\begin{figure}[H]
    \centering
    \includegraphics[width=\linewidth]{graph3_parameter_adaptation.png}
    \caption{Parameter adaptation across functions and dimensions. The meta-GA discovers distinct strategies (e.g., varying mutation rates and population sizes) tailored to each specific landscape.}
    \label{fig:param_adaptation}
\end{figure}

\subsection{Discussion}

Overall, meta-optimization confirms that the Homework~2 GA was already
well tuned for easier benchmarks (De~Jong~1 and low-dimensional Rastrigin),
since the meta-GA essentially recovers equivalent performance
(compare Tables~\ref{tab:baseline} and \ref{tab:comparison}).
However, for more complex and deceptive landscapes, the meta-optimized GA
achieves systematic and sometimes substantial gains.

Examining the parameter configurations in Table~\ref{tab:opt-params},
several interesting patterns emerge.
On Schwefel, the meta-GA tends to select:
\begin{itemize}
  \item larger populations in low and moderate dimensions (5D, 10D),
        combined with relatively high mutation rates, which encourage
        broad exploration;
  \item moderate populations and small mutation in 30D, but with a
        low $\eta_m$ (allowing occasional larger jumps), which together
        balance exploration and exploitation near promising regions.
\end{itemize}

For Michalewicz, the meta-GA prefers:
\begin{itemize}
  \item moderate-to-large populations (around 500--770 individuals);
  \item mutation rates around $0.09$ in 5D and 10D, dropping to
        about $0.013$ in 30D;
  \item relatively large distribution indices, especially in 30D,
        favouring small, precise steps once the GA has located good
        basins of attraction.
\end{itemize}

These patterns are consistent with the intuition that more complicated and
rugged landscapes benefit from more carefully tuned parameter combinations,
and that the optimal values can differ significantly between functions and
dimensions.

\section{Influence of parameters and insights from meta-optimization}

Homework~2 already discussed qualitatively the influence of population size,
number of generations, crossover and mutation rates, and distribution
indices.
This extended study complements that analysis by providing an automated way to
select these parameters and by examining the resulting configurations.

\paragraph{Population size vs. dimensionality.}
For De~Jong~1 and Rastrigin, the meta-GA tends to increase population size
with dimensionality, up to almost 1000 individuals in 30D
(Table~\ref{tab:opt-params}).
This is natural: in high dimensions a larger population helps cover the
search space and maintain diversity.
In contrast, for Schwefel the meta-GA settles on a population of around 560
individuals in 30D, significantly smaller than the 1000 used in Homework~2.
This suggests that for highly deceptive functions, a very large population is
not always necessary if the other parameters are tuned appropriately.

\paragraph{Crossover vs. mutation.}
In all cases, crossover probabilities remain high (typically above 0.75),
confirming the central role of recombination in building good solutions.
Mutation rates, however, vary widely across problems and dimensions.
Schwefel 10D uses a very high mutation rate ($p_m \approx 0.35$), which
helps escape numerous local minima, whereas De~Jong~1 and Rastrigin 30D use
much lower mutation rates (around $0.02$), emphasizing exploitation once the
global basin is found.

\paragraph{Distribution indices.}
The meta-GA often chooses large values for $\eta_c$ and $\eta_m$,
especially in high dimensions (e.g.\ Michalewicz 30D),
indicating a preference for small, fine-grained steps in the later stages of
the search.
On Schwefel, lower $\eta_m$ in 5D and 10D allows more exploratory mutations,
which appears beneficial for escaping deceptive traps.

\section{Conclusions and future work}

This extended study introduced a meta-optimization layer on top of the real-coded
GA developed in Homework~2.
By letting a meta-GA tune the main GA parameters on a per-problem and
per-dimension basis, we obtained a more systematic understanding of how
these parameters influence performance, and we achieved measurable
improvements on the harder benchmark functions.

The main conclusions are:
\begin{enumerate}
  \item On simple unimodal or already-solved multimodal functions
        (De~Jong~1 and Rastrigin), the Homework~2 GA is virtually
        optimal; meta-optimization confirms its robustness and leaves
        performance unchanged at machine precision
        (Tables~\ref{tab:baseline} and \ref{tab:comparison}).
  \item On deceptive and rugged landscapes (Schwefel, Michalewicz),
        meta-optimization yields consistent improvements in mean fitness,
        up to about $13\%$ on Schwefel 30D, while often reducing
        variance across runs (Table~\ref{tab:comparison}).
  \item The optimized parameter configurations reveal clear patterns
        (Table~\ref{tab:opt-params}):
        population size and mutation rate need to be adapted to both
        the function and the dimensionality, and optimal distribution
        indices can differ substantially from generic textbook values.
\end{enumerate}

Future work could explore:
\begin{itemize}
  \item fully \emph{self-adaptive} GAs, where parameters are encoded in
        the individuals and evolve together with the solutions;
  \item multi-objective meta-optimization that trades off solution quality
        and computational cost (e.g.\ number of function evaluations);
  \item hybrid metaheuristics that combine the meta-optimized GA with
        local search (memetic algorithms) for further refinement of
        promising individuals.
\end{itemize}

\section*{References}

\begin{enumerate}
  \item[{[1]}] Course notes and homework descriptions for the Genetic
               Algorithms course (UAIC).
  \item[{[2]}] D.\,E. Goldberg, \emph{Genetic Algorithms in Search,
               Optimization, and Machine Learning}, Addison-Wesley, 1989.
  \item[{[3]}] K. Deb, \emph{Multi-Objective Optimization Using Evolutionary
               Algorithms}, Wiley, 2001.
  \item[{[4]}] K. Deb and R. B. Agrawal, ``Simulated binary crossover for
               continuous search space,'' \emph{Complex Systems}, vol.~9,
               pp.~115--148, 1995.
  \item[{[5]}] Wikipedia contributors, ``Test functions for optimization''
               and references therein, accessed 2025.
\end{enumerate}

\end{document}